% Coursepaper Essay — coursepaper.cls
%
% A complete course-paper example with the class's course metadata and
% double-spaced final layout.  The document is intentionally self-contained:
% the coursepaper class is supplied by TeX Live and is not bundled here.
\documentclass[final]{coursepaper}
\usepackage[T1]{fontenc}
\usepackage[utf8]{inputenc}
\usepackage{booktabs}
\usepackage{hyperref}

\title{Designing Trustworthy Queueing Systems}
\author{Jordan Ellis}
\date{29 August 2026}
\college{Northbridge College of Applied Science}
\coursename{Operations and Service Design}
\coursenumber{OSD 302}
\coursesection{Section 02}
\studentnumber{NB-20418}
\instructor{Dr. Mira Shah}

\begin{document}
\maketitle

\begin{abstract}
This essay examines how a small service team can reduce waiting time without
optimising one queue at the expense of the whole customer journey.  A simple
queueing model is paired with an operational policy: measure arrival and
service variability, publish a transparent priority rule, and review the
result against equity and reliability measures.  The example shows why a
slightly lower average utilisation can produce a substantially more dependable
service.
\end{abstract}

\section{The operational question}

Queueing decisions are often framed as a choice between adding capacity and
asking customers to wait.  That framing is incomplete.  A clinic, repair
desk, or public help line is a system of linked queues, hand-offs, and
decisions.  Improving one station can simply move congestion to the next one.
The useful question is therefore not ``How fast can this station work?'' but
``What experience does the complete service promise?''

\section{A small model}

Consider a single-server desk with a mean arrival rate of
$\lambda=18$ requests per hour and a mean service rate of
$\mu=24$ requests per hour.  The utilisation is
\[
  \rho = \frac{\lambda}{\mu}=0.75.
\]
In an idealised M/M/1 model, the expected number of requests in the system is
$L=\rho/(1-\rho)=3$, and Little's law gives an expected time in the system of
$W=L/\lambda=10$ minutes.  These figures are not a forecast: real arrivals
come in bursts, service times are not memoryless, and a second specialist may
be required.  They are a baseline that makes the effect of variability
visible.

\begin{table}[ht]
  \centering
  \caption{Illustrative operating points for the same desk.}
  \begin{tabular}{lrrr}
    \toprule
    Scenario & Arrival rate & Service rate & Utilisation \\
    \midrule
    Quiet planning hour & 12/h & 24/h & 0.50 \\
    Typical hour        & 18/h & 24/h & 0.75 \\
    Busy hour            & 21/h & 24/h & 0.875 \\
    \bottomrule
  \end{tabular}
\end{table}

\section{Policy and measurement}

The desk should reserve a small portion of capacity for urgent cases while
keeping the rule visible to everyone.  A daily dashboard can record arrival
rate, service-time percentiles, abandonment, and the fraction of requests
resolved without a second hand-off.  The service owner should inspect the
90th-percentile wait alongside the mean: the mean can improve while a small
group experiences an unacceptable delay.

Three safeguards make the policy auditable.  First, record the reason for each
priority decision.  Second, publish a weekly sample of anonymised outcomes.
Third, change one operational parameter at a time so that a claimed
improvement has a plausible cause.  These practices turn a mathematical
model into a learning loop rather than a promise of exact prediction.

\section{Conclusion}

The model makes a simple point: as utilisation approaches one, variability is
converted into waiting.  A trustworthy queueing policy therefore values
headroom, clear priority rules, and distributional measures of service quality.
The result is not merely a shorter average queue; it is a service whose
performance can be explained and improved.

\begin{thebibliography}{9}
\bibitem{little1961}
J. D. C. Little, ``A proof for the queuing formula: L = $\lambda$ W,''
\emph{Operations Research}, vol.~9, no.~3, pp.~383--387, 1961.

\bibitem{hopp2011}
W. J. Hopp and M. L. Spearman, \emph{Factory Physics}, 3rd ed.
Waveland Press, 2011.
\end{thebibliography}

\end{document}
